\documentclass[11pt]{article}
\usepackage{mathblatt}
\makeatletter
\setlength{\mbpfbe}{0mm}% keine Punkte (Bauregel 6.4)
\makeatother
% Rechenraum als Karo (Bauregel 6.9): n Rechenschritte = n mal 10 mm
\newcommand{\karo}[1]{\par\smallskip\noindent\begin{tikzpicture}
  \clip (0,0) rectangle ({\linewidth-1mm},{#1*10mm});
  \draw[black!16,line width=0.3pt,step=5mm] (0,0) grid ({\linewidth},{#1*10mm});
  \end{tikzpicture}\par}
\newcommand{\linie}[1]{\,{\color{black!45}\rule[-1pt]{#1}{0.4pt}}\,}
\newcommand{\pkt}{\textperiodcentered{}}
\begin{document}

\pfheftstil
\pflaufkopf[P10 \pkt{} mehr]{Vierecke}{PG30-2}\pfseitevon
\pfheftkopf{Vierecke}{mehr}
\pfstufenzeile{Symmetrieachsen zählen \pkt{} Eigenschaft erkennen}
% 2021-OS-B1j-m1 (Schwere I, rang 4)
\begin{pfaufg}{nach~P10\\’21~\pkt{}~1j}{1.}{}
Wie viele Symmetrieachsen hat ein Quadrat? Kreuze an:
\par\smallskip\noindent{\small\itshape Skizze: Viereck (Quadrat): grau gefüllt}\par
\par\smallskip\noindent\pfkreuz{A}{0}\pfkreuz{B}{1}\pfkreuz{C}{2}\pfkreuz{D}{3}
\fusshilfe{\mbox{1:~4 Symmetrieachsen (Option „4“)}}
\end{pfaufg}
% 2016-OS-B1e-m1 (Schwere I, rang 5)
\begin{pfaufg}{nach~P10\\’16~\pkt{}~1e}{2.}{}
Kreuze die Aussage an, die für jedes Parallelogramm gilt:
\par\smallskip\noindent\pfkreuz{A}{Je zwei benachbarte Winkel sind gleich.}\pfkreuz{B}{Je zwei gegenüberliegende Winkel sind gleich.}\pfkreuz{C}{Alle vier Winkel sind gleich.}
\fusshilfe{\mbox{2:~B}}
\end{pfaufg}
\par\vfill\noindent{\scriptsize\color{mbgrau}Originale: P10 2021 \pkt{} 1j, P10 2016 \pkt{} 1e\par}
\end{document}